\section{Experiments}
\label{sec:exp}
%% Camera-ready rewrite. Submitted text: sec/4_experiments.tex. All numbers come from
%% textfix_numbers.json; tables are generated by make_camera_ready_tables.py.

\subsection{Datasets and Evaluation Protocol}
\label{sec:datasets}

We train on DUTS-TR~\cite{wang2017duts} ($10{,}553$ images) and evaluate on
five standard RGB SOD benchmarks: DUTS-TE ($5{,}019$)~\cite{wang2017duts},
DUT-OMRON ($5{,}168$)~\cite{yang2013duto}, HKU-IS ($4{,}447$)~\cite{li2015hkuis},
ECSSD ($1{,}000$)~\cite{yan2013ecssd}, and PASCAL-S ($850$)~\cite{li2014pascals},
plus SIP~\cite{fan2019sip}. \rev{SIP is not used for training, so we evaluate on all
of its $929$ images, as the RGB-D methods we compare against do. We report
S-measure ($S_m$)~\cite{fan2017smeasure}, mean E-measure
($E_m$)~\cite{fan2018emeasure}, and mean absolute error ($\mathcal{M}$), computed on
the continuous saliency map with the PySODMetrics implementation, and the
F-measure ($F_\beta$, $\beta^2{=}0.3$) at a fixed threshold of $0.5$, averaged over
images. The threshold is not tuned on any dataset.} No test-time augmentation or
post-processing is applied. \rev{We express differences in $S_m$ in points
($0.01 = 1$ point).}

\subsection{Implementation Details}
\label{sec:impl}

\begin{revpar}
We use CLIP ViT-L/14 at $336{\times}336$ and SAM ViT-B, with $N{=}4$ ESC blocks,
$N_o{=}5$ pseudo-prompt regions and $K{=}10$ prompts. The CLIP text encoder is
frozen. The CLIP visual encoder is fine-tuned: its attention layers and the reused
SAM blocks use a learning rate of $2{\times}10^{-6}$, and its remaining layers and
all new modules use $2{\times}10^{-5}$. We train for $30$ epochs with AdamW (weight
decay $10^{-4}$), batch size $8$, a 5-epoch linear warm-up and cosine decay with
$10{\times}$ step-downs at $50\%$ and $75\%$ of training, and standard geometric
and photometric augmentation. The checkpoint with the best training Dice is kept,
so no test benchmark is used for model selection. A run takes about $5$\,h on one
A100 GPU. Comparisons with prior work use seed 42; the study reports three seeds
(42, 1, 2) for the main configurations.
\end{revpar}

\subsection{Comparison with State-of-the-Art Methods}
\label{sec:sota}

% AUTO-GENERATED by make_camera_ready_tables.py from textfix_numbers.json
\begin{table*}[t]
  \centering
  \caption{Comparison with RGB SOD methods. Best $S_m$ and $\mathcal{M}$ per column in
    \textbf{bold}. \rev{Competitor values are taken from published tables (see text);
    their $F_\beta$ is the maximum F-measure, or a variant not stated in the source
    ($^\dagger$), while ours is $F_\beta$ at a fixed threshold of $0.5$, so the
    $F_\beta$ column is not directly comparable and is not bolded. \methodname{}:
    Baseline+MP+AF, CLIP ViT-L/14 at $336^2$, seed 42; three-seed means are in
    \cref{tab:prompt_side}.}}
  \label{tab:sota}
  \setlength{\tabcolsep}{2.6pt}
  \scriptsize
  \begin{tabular}{l ccc ccc ccc ccc ccc}
    \toprule
    & \multicolumn{3}{c}{DUTS-TE} & \multicolumn{3}{c}{DUT-OMRON}
    & \multicolumn{3}{c}{ECSSD} & \multicolumn{3}{c}{HKU-IS} & \multicolumn{3}{c}{\rev{PASCAL-S}} \\
    \cmidrule(lr){2-4}\cmidrule(lr){5-7}\cmidrule(lr){8-10}\cmidrule(lr){11-13}\cmidrule(lr){14-16}
    Method & $S_m$ & $F_\beta$ & $\mathcal{M}$ & $S_m$ & $F_\beta$ & $\mathcal{M}$
           & $S_m$ & $F_\beta$ & $\mathcal{M}$ & $S_m$ & $F_\beta$ & $\mathcal{M}$
           & $S_m$ & $F_\beta$ & $\mathcal{M}$ \\
    \midrule
    PiCANet~\cite{picanet2018} & .869 & .860 & .051 & .832 & .803 & .065 & .917 & .935 & .046 & .904 & .918 & .043 & \rev{.852} & \rev{.868} & \rev{.078} \\
    BASNet~\cite{basnet2019} & \rev{.866} & .860 & .048 & .836 & .805 & .056 & .916 & .942 & .037 & .909 & .928 & .032 & \rev{.834} & \rev{.860} & \rev{.079} \\
    F3Net~\cite{f3net2020} & .888 & .891 & .035 & .838 & .813 & .053 & .924 & .945 & .033 & .917 & .937 & .028 & \rev{.857} & \rev{.882} & \rev{.064} \\
    LDF~\cite{ldf2020} & .892 & .898 & .034 & \rev{.838} & .820 & \rev{.051} & \rev{.924} & .950 & .034 & .919 & .939 & .027 & \rev{.859} & \rev{.887} & \rev{.062} \\
    \chk{CSF-R2}~\cite{csnet2020} & \rev{.890} & \rev{.890} & \rev{.037} & \rev{.838} & \rev{.815} & .055 & \rev{.930} & \rev{.950} & \rev{.033} & .921 & .935 & .030 & \rev{.862} & \rev{.886} & \rev{.069} \\
    VST~\cite{vst2021} & .896 & .890 & .037 & .850 & .825 & .058 & .932 & .951 & .033 & .928 & .942 & .029 & \rev{.871} & \rev{.890} & \rev{.062} \\
    PSGLoss~\cite{psgloss2021} & \rev{.883} & \rev{.886} & .036 & \rev{.831} & \rev{.811} & .052 & \rev{.925} & \rev{.949} & \rev{.031} & .919 & .938 & .027 & \rev{.858} & \rev{.886} & \rev{.063} \\
    \rev{RCSB~\cite{rcsb2022}} & \rev{.878} & \rev{.889} & \rev{.035} & \rev{.820} & \rev{.810} & \rev{.045} & \rev{.921} & \rev{.944} & \rev{.033} & \rev{.918} & \rev{.938} & \rev{.027} & \rev{.857} & \rev{.886} & \rev{.061} \\
    EDN~\cite{edn2022} & .892 & .895 & .035 & \rev{.846} & .828 & \rev{.048} & .927 & .951 & .032 & .924 & .941 & .026 & \rev{.860} & \rev{.891} & \rev{.065} \\
    MENet~\cite{menet2023} & .905 & \rev{.913} & .028 & .850 & .834 & .045 & .928 & \rev{.955} & .031 & .927 & .948 & .023 & \rev{.868} & \rev{.901} & \rev{.057} \\
    DC-Net$^\dagger$~\cite{sod_dcnet} & \rev{.896} & \rev{.829} & \rev{.035} & \rev{.849} & \rev{.774} & \rev{.053} & \rev{.924} & \rev{.882} & \rev{.034} & .924 & .910 & .027 & \rev{.857} & \rev{.769} & \rev{.066} \\
    ADMNet$^\dagger$~\cite{sod_admnet} & \rev{.835} & \rev{.753} & \rev{.057} & \rev{.815} & \rev{.722} & \rev{.063} & \rev{.889} & \rev{.870} & \rev{.051} & .889 & .871 & .040 & \rev{.798} & \rev{.731} & \rev{.094} \\
    BDFGNet$^\dagger$~\cite{sod_bdfgnet} & \rev{.884} & \rev{.794} & \rev{.049} & \rev{.826} & \rev{.723} & \rev{.058} & \rev{.908} & \rev{.891} & \rev{.042} & .909 & .882 & .042 & \rev{.844} & \rev{.791} & \rev{.069} \\
    STNet$^\dagger$~\cite{sod_stnet} & \rev{.860} & \rev{.813} & \rev{.044} & \rev{.823} & \rev{.752} & \rev{.056} & \rev{.909} & \rev{.914} & \rev{.040} & .902 & .898 & .033 & \rev{.840} & \rev{.821} & \rev{.072} \\
    \rev{MDSAM~\cite{mdsam2024} (SAM, 512$^2$)} & \rev{.920} & \rev{.937} & \rev{\textbf{.024}} & \rev{.878} & \rev{.887} & \rev{\textbf{.039}} & \rev{.948} & \rev{.974} & \rev{\textbf{.021}} & \rev{\textbf{.941}} & \rev{.963} & \rev{\textbf{.019}} & \rev{.882} & \rev{.907} & \rev{.052} \\
    \midrule
    \methodname{} (ours) & \rev{\textbf{.929}} & \rev{.913} & \rev{.025} & \rev{\textbf{.884}} & \rev{.840} & \rev{.045} & \rev{\textbf{.949}} & \rev{.955} & \rev{.024} & \rev{.940} & \rev{.941} & \rev{.025} & \rev{\textbf{.896}} & \rev{.878} & \rev{\textbf{.050}} \\
    \bottomrule
  \end{tabular}
\end{table*}


\begin{revpar}
\cref{tab:sota} compares \methodname{} with CNN-, transformer- and SAM-based RGB
SOD methods. Competitor values are taken from published comparison tables: the
classical methods from Dulam and Kambhamettu~\cite{dulam2025soddcnet}, DC-Net,
ADMNet, BDFGNet and STNet from the STNet paper~\cite{sod_stnet}, and MDSAM from its
own paper~\cite{mdsam2024}; we corrected submitted values that did not match these
sources. \methodname{} obtains the highest $S_m$ on four of the five benchmarks:
DUTS-TE ($0.929$ vs.\ $0.920$ for MDSAM), DUT-OMRON ($0.884$ vs.\ $0.878$), ECSSD
($0.949$ vs.\ $0.948$) and PASCAL-S ($0.896$ vs.\ $0.882$). MDSAM, which adapts SAM
with multi-scale adapters at $512{\times}512$, is slightly higher on HKU-IS
($0.941$ vs.\ $0.940$) and has the lowest MAE on four benchmarks. On ECSSD and
HKU-IS the margins are within our seed-to-seed variation, so we regard the two
methods as on par there.
\end{revpar}

\noindent\textbf{RGB-D comparison on SIP.}
% AUTO-GENERATED by make_camera_ready_tables.py from textfix_numbers.json
\begin{table}[t]
  \centering
  \caption{SIP (all 929 images): \methodname{} (RGB only, trained on DUTS-TR) vs.\
    RGB-D methods (RGB{+}depth, trained on RGB-D data). \rev{RGB-D values are from
    VST++~\cite{vstpp2024} (max $F_\beta$, max $E_m$); \methodname{} values are $F_\beta$ at threshold
    $0.5$ and mean $E_m$ (our max $F_\beta$/max $E_m$: .947/.969).
    Three seeds: $S_m$ .935$\pm$0.004, $\mathcal{M}$ .028$\pm$0.002.
    Best per column in \textbf{bold}. Purple: changed from the submitted version; please verify.}}
  \label{tab:sip_rgbd}
  \setlength{\tabcolsep}{4pt}
  \footnotesize
  \begin{tabular}{llcccc}
    \toprule
    Method & Input & $S_m\!\uparrow$ & $F_\beta\!\uparrow$ & $E_m\!\uparrow$ & $\mathcal{M}\!\downarrow$ \\
    \midrule
    \rev{S2MA~\cite{sod_s2ma}} & RGB-D & \rev{.872} & \rev{.877} & \rev{.919} & \rev{.058} \\
    \rev{BBS-Net~\cite{sod_bbsnet}} & RGB-D & \rev{.879} & \rev{.884} & \rev{.922} & \rev{.055} \\
    \rev{JL-DCF~\cite{sod_jldcf}} & RGB-D & \rev{.885} & \rev{.894} & \rev{.931} & \rev{.049} \\
    \rev{SPNet~\cite{sod_spnet}} & RGB-D & \rev{.894} & \rev{.904} & \rev{.933} & \rev{.043} \\
    \rev{SPSN~\cite{sod_spsn}} & RGB-D & \rev{.892} & \rev{.900} & \rev{.936} & \rev{.043} \\
    VST~\cite{vst2021} & RGB-D & .904 & \chk{.915} & .944 & .040 \\
    EBMGSOD~\cite{sod_ebm} & RGB-D & .902 & \chk{.918} & \chk{.944} & .037 \\
    \midrule
    \methodname{} & RGB & \rev{\textbf{.936}} & \rev{\textbf{.937}} & \rev{\textbf{.958}} & \rev{\textbf{.027}} \\
    \bottomrule
  \end{tabular}
\end{table}

\rev{SIP~\cite{fan2019sip} contains people in real scenes and is usually addressed
with RGB-D methods trained on RGB-D data. Using RGB input only and training on
DUTS-TR, \methodname{} reaches $S_m{=}0.936$ and $\mathcal{M}{=}0.027$
(\cref{tab:sip_rgbd}), above VST~\cite{vst2021} ($0.904$, $0.040$) and
EBMGSOD~\cite{sod_ebm} ($0.902$, $0.037$). The training data differ, so this shows
that a strong RGB foundation model is competitive on SIP; it does not isolate the
contribution of depth.}

\subsection{Study of Prompt-Side Adaptations}
\label{sec:study_prompt}

% AUTO-GENERATED by make_camera_ready_tables.py from textfix_numbers.json
\begin{table*}[t]
  \centering
  \begin{revpar}
  \caption{Prompt-side adaptations, all with the simple decoder and BCE loss. Runs: number of seeds (42, 1, 2); multi-seed entries are mean$\pm$std. $\overline{\mathcal{M}}$: MAE averaged over the six benchmarks. MP: 10-prompt bank with max aggregation; AF: adaptive fallback; VLF-$K$: fusion block attends over the $K$ bank prompts and the caption. \textbf{Ours} = MP+AF.}
  \label{tab:prompt_side}
  \setlength{\tabcolsep}{3.2pt}
  \footnotesize
  \begin{tabular}{lc cccccc c}
    \toprule
    & & \multicolumn{6}{c}{$S_m\uparrow$} & \\
    \cmidrule(lr){3-8}
    Configuration & Runs & DUTS-TE & DUT-OMRON & ECSSD & HKU-IS & PASCAL-S & SIP & $\overline{\mathcal{M}}\downarrow$ \\
    \midrule
    Baseline (single prompt) & 3 & .926{\scriptsize$\pm$.000} & .879{\scriptsize$\pm$.002} & .950{\scriptsize$\pm$.001} & .940{\scriptsize$\pm$.001} & .897{\scriptsize$\pm$.002} & .933{\scriptsize$\pm$.003} & .034 \\
    + AF & 1 & .925 & .880 & .951 & .940 & .898 & .931 & .034 \\
    + MP & 1 & .926 & .885 & .951 & .940 & .898 & .931 & .033 \\
    \textbf{+ MP + AF (ours)} & 3 & .928{\scriptsize$\pm$.001} & .883{\scriptsize$\pm$.002} & .949{\scriptsize$\pm$.000} & .940{\scriptsize$\pm$.000} & .897{\scriptsize$\pm$.001} & .935{\scriptsize$\pm$.004} & .033 \\
    \midrule
    Ours, $K{=}1$ & 1 & .927 & .880 & .952 & .941 & .897 & .938 & .033 \\
    Ours, $K{=}5$ & 1 & .927 & .881 & .950 & .940 & .899 & .934 & .033 \\
    Ours, random wording & 1 & .923 & .876 & .952 & .940 & .896 & .935 & .035 \\
    Ours, diverse wording & 1 & .926 & .883 & .951 & .941 & .898 & .933 & .033 \\
    \midrule
    Ours + VLF-$K$ & 3 & .925{\scriptsize$\pm$.003} & .877{\scriptsize$\pm$.005} & .949{\scriptsize$\pm$.000} & .940{\scriptsize$\pm$.001} & .897{\scriptsize$\pm$.001} & .933{\scriptsize$\pm$.002} & .035 \\
    Ours + VLF-$K$, no in-block residual & 1 & .926 & .879 & .948 & .939 & .897 & .934 & .033 \\
    \bottomrule
  \end{tabular}
  \end{revpar}
\end{table*}


\begin{revpar}
\cref{tab:prompt_side} starts from the single-prompt baseline and varies only the
text side, keeping the simple decoder and BCE loss. Seed-to-seed standard
deviations of the three-seed rows are $0.0$--$0.5$ points.

\noindent\textbf{Prompt bank and adaptive fallback.} Combining the 10-prompt bank
with the adaptive fallback (ours) improves $S_m$ over the baseline by $0.2$ points
on DUTS-TE and $0.3$ points on DUT-OMRON, larger than the seed standard deviations,
and is on par
on the other four benchmarks. Neither component explains the gain alone: the
fallback without a bank leaves the baseline unchanged, and the bank without the
fallback is higher on DUT-OMRON but lower on SIP.

\noindent\textbf{Bank size and wording.} With $K{=}1$ or $K{=}5$, or with randomly
sampled or deliberately diverse wording, $S_m$ stays within $0.4$ points of ours on
every benchmark, except random wording on DUT-OMRON ($-0.6$). For reference, two single-prompt runs that differ only in their
code path ($+$AF and ours with $K{=}1$) differ by $0.7$ points on SIP, so
single-run differences of this size are not meaningful. With a fine-tuned
ViT-L/14, the exact prompts matter little.

\noindent\textbf{Multi-token fusion.} Letting the VLF block attend over the bank
prompts (VLF-$K$) makes its output depend on its inputs, but it does not help: over
three seeds it is $0.3$ and $0.5$ points lower on DUTS-TE and DUT-OMRON and on par
elsewhere. The single-key behaviour of \cref{sec:sam} is therefore not what limits
accuracy.
\end{revpar}

\subsection{Study of Architecture and Training Adaptations}
\label{sec:study_arch}

% AUTO-GENERATED by make_camera_ready_tables.py from textfix_numbers.json
\begin{table*}[t]
  \centering
  \begin{revpar}
  \caption{Architecture and training adaptations. ``Corr-aware'': \methodname{} + bidirectional SAM attention, ASPP/FPN decoder that reads the correlation, compound loss and deep supervision. Residual: $C_{\textit{final}}{=}C_N{+}C_0$ (\cref{eq:residual}). Frozen: CLIP visual encoder and SAM blocks frozen, new modules trained. Columns as in \cref{tab:prompt_side}.}
  \label{tab:arch_side}
  \setlength{\tabcolsep}{3.2pt}
  \footnotesize
  \begin{tabular}{lc cccccc c}
    \toprule
    & & \multicolumn{6}{c}{$S_m\uparrow$} & \\
    \cmidrule(lr){3-8}
    Configuration & Runs & DUTS-TE & DUT-OMRON & ECSSD & HKU-IS & PASCAL-S & SIP & $\overline{\mathcal{M}}\downarrow$ \\
    \midrule
    \methodname{} (simple decoder, BCE) & 3 & \textbf{.928}{\scriptsize$\pm$.001} & \textbf{.883}{\scriptsize$\pm$.002} & .949{\scriptsize$\pm$.000} & \textbf{.940}{\scriptsize$\pm$.000} & \textbf{.897}{\scriptsize$\pm$.001} & \textbf{.935}{\scriptsize$\pm$.004} & \textbf{.033} \\
    \methodname{} + BiSAM & 1 & .926 & .880 & \textbf{.951} & \textbf{.940} & .895 & .932 & \textbf{.033} \\
    \midrule
    Corr-aware, max, residual on & 3 & .926{\scriptsize$\pm$.001} & .882{\scriptsize$\pm$.001} & .948{\scriptsize$\pm$.001} & .937{\scriptsize$\pm$.002} & .896{\scriptsize$\pm$.002} & .934{\scriptsize$\pm$.005} & .035 \\
    Corr-aware, max, residual off & 3 & .925{\scriptsize$\pm$.003} & .878{\scriptsize$\pm$.004} & .949{\scriptsize$\pm$.002} & .938{\scriptsize$\pm$.001} & .896{\scriptsize$\pm$.001} & .932{\scriptsize$\pm$.002} & .036 \\
    Corr-aware, attention, residual on & 3 & .924{\scriptsize$\pm$.000} & .878{\scriptsize$\pm$.002} & .949{\scriptsize$\pm$.001} & .938{\scriptsize$\pm$.001} & .895{\scriptsize$\pm$.002} & .934{\scriptsize$\pm$.000} & .035 \\
    Corr-aware, attention, residual off & 3 & .926{\scriptsize$\pm$.001} & .882{\scriptsize$\pm$.001} & .950{\scriptsize$\pm$.001} & .938{\scriptsize$\pm$.000} & .895{\scriptsize$\pm$.003} & .933{\scriptsize$\pm$.003} & .034 \\
    Corr-aware, max, residual on, + point loss & 3 & .921{\scriptsize$\pm$.002} & .875{\scriptsize$\pm$.002} & .945{\scriptsize$\pm$.002} & .935{\scriptsize$\pm$.002} & .892{\scriptsize$\pm$.002} & .931{\scriptsize$\pm$.002} & .037 \\
    \midrule
    \methodname{}, frozen encoders & 3 & .903{\scriptsize$\pm$.001} & .839{\scriptsize$\pm$.001} & .932{\scriptsize$\pm$.001} & .924{\scriptsize$\pm$.001} & .889{\scriptsize$\pm$.001} & .916{\scriptsize$\pm$.002} & .042 \\
    \bottomrule
  \end{tabular}
  \end{revpar}
\end{table*}


\begin{revpar}
\noindent\textbf{Decoder and loss.} Adding BiSAM to ours, or switching to the
correlation-aware variant (BiSAM, ASPP/FPN decoder, compound loss and deep
supervision), gives $S_m$ within $0.3$ points of ours on every benchmark (all
correlation-aware settings stay within $0.5$) and a slightly higher MAE
(\cref{tab:arch_side}). The simple decoder with BCE is therefore sufficient. The
auxiliary point-classification loss lowers $S_m$ of the correlation-aware variant
on all six benchmarks, by $0.4$ points on DUTS-TE and $0.7$ on DUT-OMRON.

\noindent\textbf{Prompt residual and aggregation.} In the correlation-aware variant
the prompt residual is the only path from the loss to the prompt branch through the
correlation: over 50 DUTS-TR batches, the mean gradient norm reaching $W_v$ is
$0.75$ with the residual and exactly $0$ without it. Restoring this gradient does
not translate into accuracy. Turning the residual off changes $S_m$ by at most
$0.4$ points, with opposite signs for max-pooling and attention aggregation, and
attention is not better than max-pooling. We conclude that neither choice has a
consistent effect.

\noindent\textbf{Encoder adaptation.} Freezing the CLIP visual encoder and the SAM
blocks, while still training all new modules, lowers $S_m$ by $2.4$, $4.4$, $1.8$,
$1.6$, $0.8$ and $1.8$ points on the six benchmarks and raises the mean MAE from
$0.033$ to $0.042$. This is by far the largest effect in our study.
\end{revpar}

\subsection{Discussion}
\label{sec:discussion}

\begin{revpar}
Three observations summarise the study. \emph{(1) Adapt the encoder first.}
Fine-tuning the CLIP visual encoder contributes several times more than any
prompt-side or architectural choice; for SOD, adapting CLIP's features to
saliency matters more than how the text side is designed. \emph{(2) Keep the
pipeline simple.} The prompt bank with adaptive fallback gives small gains on the
hardest benchmarks, while the correlation-aware decoder, compound loss, BiSAM,
attention aggregation, the prompt residual and multi-token fusion add complexity
without a consistent benefit, and the point loss hurts. \emph{(3) Test that each
branch is used.} A single-key text attention, and an implementation that skipped
the text transformer, both still trained to competitive accuracy; simple
input-sensitivity and gradient checks expose such silent failures, and we
recommend them when porting multi-class foundation-model pipelines to binary tasks.
\end{revpar}
